\setcounter{section}{1}\setcounter{subsection}{0}
% Original source lines 90--204
\subsection{Notations and conventions}
In this article, $M$ denotes a $d$-dimensional (with $d \geq 3$) smooth manifold endowed with a metric tensor $g$ or conformal class of metrics $\cc$, both of which will be positive definite. Furthermore, $\Lambda$~denotes a $(d-k)$-dimensional smooth submanifold of $M$, and $0 < k < d$. Throughout the entirety of this article, we work only locally and so will assume that trivializations are always available for any section of any bundle. In the Riemannian setting $(M,g)$, we denote the Levi-Civita connection by $\nabla$. Our convention for the curvature tensor is given~by
\[
R(x,y)z = (\nabla_x \nabla_y - \nabla_y \nabla_x)z - \nabla_{[x,y]} z .
\]
Here, $x$, $y$, $z$ are vector fields on $M$ and $[x,y]$ is the Lie bracket of $x$ and $y$. In general, for any bundle $\mathcal{V}M$ over $M$, we will denote sections of this bundle by $\Gamma(\mathcal{V} M)$. That is, $x,y,z \in \Gamma(TM)$. When necessary, we will refer to arbitrary tensor products of a bundle (and its dual) using capital Greek letter, so that, for example, $T^\Phi M$ is some unspecified tensor product of $TM$ and~$T^* M$.

For ease of calculation, we will use Penrose's abstract index notation, labelling such indices with letters from the beginning of the Latin alphabet $a$, $b$, $c$, etc., so that, for example, the~Riemann curvature tensor may be expressed as
\[
R_{ab}{}^c{}_d x^d = [\nabla_a, \nabla_b] x^c .
\]
Using the metric and its inverse, indices can be raised and lowered, as usual; the Einstein summation convention is used to indicate contraction. (Sometimes, when the notation is unambiguous, we will use dot product notation to represent contraction; for example, when $\omega$ is some 1-form, we will sometimes write $\nabla^a \omega_a =: \nabla \cdot \omega$.) In this language, we have $\mathrm{Ric}_{ab} := R_{ca}{}^c{}_b$ and its corresponding scalar is $\mathrm{Sc} := \mathrm{Ric}_a^a$. When $d > 2$, a particular trace-correction of the Ricci tensor, known as the Schouten tensor, arises frequently in conformal considerations,
\[
P_{ab} := \frac{1}{d-2} \left(\mathrm{Ric}_{ab} - \frac{1}{2(d-1)} \mathrm{Sc}\, g_{ab} \right) .
\]
The trace of this tensor is denoted by $J := P_a^a$. Written in terms of the Riemann tensor and the Schouten tensor, we may express the Weyl tensor as
\[
W_{abcd} = R_{abcd} - g_{ac} P_{bd} + g_{bc} P_{ad} + g_{ad} P_{bc} - g_{bd} P_{ac} .
\]
In dimensions $d > 2$, the Cotton tensor is given~by the covariant curl of the Schouten tensor,
\[
C_{abc} = \nabla_a P_{bc} - \nabla_b P_{ac} .
\]
When $d > 3$, this can be reexpressed in terms of the divergence of the Weyl tensor,
\[
C_{abc} = \frac{1}{d-3} \nabla^d W_{dcab} .
\]

In the context of submanifold embeddings, we will use Greek indices to represent tensors on normal frame bundles. Objects that are along the submanifold will be denoted by overbars~$\bar{\bullet}$, so that, for example, $\bar{\nabla}$ represents the induced Levi-Civita connection on $(\Lambda,\bar{g})$, where $\bar{g}$ is the induced metric from $\Lambda \hookrightarrow (M,g)$. In general for any index notation, we will use round brackets~$(\cdots)$ to indicate unit-normalized symmetrization and square brackets $[\cdots]$ to indicate unit-normalized antisymmetrization. Furthermore, we will use a subscript $\circ$ to indicate the tracefree part of a particular symmetric index configuration. However, because Latin and Greek indices represent different types of objects, we may sometimes write $(\alpha a b \beta)$ to specifically mean the symmetrization over $\alpha$ and $\beta$, and similarly for antisymmetrization. To be explicit, whichever type of index is closest to the opening and closing brackets is the type of index that is (anti)symmetrized over. On the other hand, when excluding an index from (anti)symmetrization is necessary, we will use vertical bars, so that $(\alpha | \beta | \gamma)$ refers to the symmetrization over~$\alpha$ and~$\gamma$. Note that Greek indices will always be subscripts, so repeated Greek indices will denote summation in the usual way.

\section[Orthonormal frames for normal bundles of Riemannian submanifold embeddings]{Orthonormal frames for normal bundles \\ of Riemannian submanifold embeddings} \label{riem-frames}

We begin with the case of smooth Riemannian submanifold embeddings $\iota \colon \Lambda^{d-k} \hookrightarrow (M,g)$ with $d \geq 3$. For notational convenience, we will often view $\Lambda$ as a subset of $M$. We now review classical submanifold geometry.

\subsection{Classical submanifold geometry}

 At a point $p \in \Lambda$, every tangent space of the bulk space along $\Lambda$ can be decomposed according~to
\begin{align} \label{TM-decomp}
T_p M = T_p \Lambda \oplus N_p \Lambda ,
\end{align}
where $N_p \Lambda$ is the normal space at $p \in \Lambda$. The normal bundle over $\Lambda$ is defined as $N \Lambda := \sqcup_{p \in \Lambda} N_p \Lambda$.

As the pullback map is well defined, we can construct the \textit{intrinsic metric}, also sometimes called the \textit{first fundamental form}, via the pullback: $\bar{g} := \iota^* g$. Note that this metric is compatible with the decomposition above, and given an orthornormal frame $\{n_{\alpha}\}$ for $N \Lambda$, we may express this invariant in index notation according to $\bar{g}_{ab} = g_{ab}|_{\Lambda} - n_{a \alpha} n_{b \alpha}$. In the notation $n^a_{\alpha}$ (and $n_{a \alpha}$), the Latin index is an abstract index whereas the Greek letter indexes some particular element of the basis -- this will always be the case. Closely related to the first fundamental form is the projection operator given~by $\bar{g}^a_b := \delta^a_b - n^a_{\alpha} n_{b \alpha}$.

With a metric in hand, one may construct the Levi-Civita connection on $\Lambda$, denoted $\bar{\nabla}$. A~classical result of submanifold geometry is that this connection is directly obtainable from the Levi-Civita connection $M$ and the projection operator,
\[
\bar{\nabla}_{\bar v} \bar{u}^a = \bar{g}^a_b \nabla_{\iota_* \bar{v}} (\iota_* \bar{u})^b ,
\]
where $\bar{v},\bar{u} \in \Gamma(T \Lambda)$ and $\iota_* \bar{v}$, $\iota_* \bar{u}$ are their respective pushforwards. Note that we will often write~$\nabla^\top$ for the pullback of the bulk Levi-Civita connection -- this way, we may use abstract index notation and we do not need to carry around arbitrary vector fields. In a similar way, we use $\top$ both as an operator and as a superscript to denote the projection to the submanifold, e.g., we will write $P^\top$ to represent the projection of the Schouten tensor of $(M,g)$ to $\Lambda$.

If we instead project $\nabla_{\iota_* \bar{v}} (\iota_* \bar{u})$ along a unit normal vector, we obtain a classical result: the~\textit{second fundamental form}. Indeed, relative to an orthonormal frame on $N \Lambda$, we define
\[
\bar{u}^a \bar{v}^b \II_{ab \alpha} := -n_{a \alpha} \nabla_{\iota_* \bar{v}} (\iota_* \bar{u})^a .
\]
A short calculation shows that this can be expressed in a simpler way according to
\[
\II_{ab \alpha} = \bar{g}_{bc} \nabla^\top_a n^c_{\alpha} \in \Gamma\bigl(\odot^2 T^* \Lambda\bigr) .
\]
It is worth pointing out that in our calculations, we view $\II_{ab \alpha}$ as a tuple of 2-tensors depending on an orthonormal frame, as per the above section space. However when an orthonormal frame is not prescribed, the second fundamental form is still well defined, taking values in the normal bundle $N \Lambda$.
Accordingly, one may express the Levi-Civita connection of $M$ restricted to $\Lambda$ via the following identity:
\[
\nabla_{\iota_* \bar{v}} (\iota_* \bar{u}) = \bar{\nabla}_{\bar{v}} \bar{u} - \II_{\alpha}(\bar{u},\bar{v}) n_{\alpha} .
\]

The first and second fundamental forms are sufficient to completely describe an embedded hypersurface in $\mathbb{E}^d$ up to Euclidean motions. This is because, for $k=1$, the normal bundle is one-dimensional and hence (up to orientation) each normal space has a canonical basis given~by $n \in N_p \Lambda$, where $g(n,n) = 1$. (Note that as we work locally for the entirety of this article, we will neglect any issues that arise from orientations.) This rigidity effectively removes all geometric structure from the normal bundle. Unfortunately, this rigidity is lost entirely when we consider higher-codimension submanifolds. Indeed, while orthonormality can certainly be retained, there is (a priori) no uniquely preferred basis for $N \Lambda$, even locally. In what follows, we will find that there is sometimes a preferred basis up to constant ${\rm SO}(k)$ rotations -- but in general, it is not that clear that any preferred basis can be constructed. We will now provide a construction for the final piece of data required to describe embedded submanifolds up to Euclidean motions and also provide notation and language that will be necessary for later considerations.

Just as one can induce a connection on the tangent bundle to $\Lambda$, one may also induce on the normal bundle a connection
\[
D \colon \ \Gamma(N \Lambda) \rightarrow \Gamma(T^* \Lambda \otimes N \Lambda)
\]
given~by
\[
D_{\bar v} \xi = {\perp} (\nabla_{\iota_* \bar{v}} \xi) ,
\]
where $\perp$ is the orthogonal projection to $N \Lambda$, $\bar{v} \in \Gamma(T \Lambda)$ is any vector field along $\Lambda$ and $\xi \in \Gamma(N \Lambda)$ is any normal vector field. Note that this connection is compatible with the metric on the normal bundle given~by $\eta = g|_{\Lambda} - \iota^* g$. Then, for any frame $\{n_{\alpha}\}$ for $N \Lambda$, the connection coefficients are defined by
\[
D_a n^b_{\beta} = \omega_{\gamma a \beta} n^b_{\gamma} .
\]
Given any orthonormal frame $\{n_{\alpha}\}_{\alpha = 1}^k$ for the normal bundle $N \Lambda$, $\omega$ is given precisely by the so-called \textit{normal fundamental forms}~\cite{Spivak} (sometimes called the H{\'a}j{\'i}{\v{c}}ek one-form in the physics literature~\cite{hajicek1975, hajicek1973}):
\[
\omega_{\alpha a \beta} = \beta_{a \alpha \beta} := n^b_{\alpha} \nabla^\top_a n_{b \beta} .
\]
Note that we use $\beta$ only for orthonormal frames.

Associated to the normal bundle is the orthonormal frame bundle $F(N \Lambda)$, which is a principal ${\rm O}(k)$-bundle. So in particular, any two orthonormal frame fields are related to one another by the action of a section $m \in C^\infty(\Lambda, {\rm O}(k))$. Furthermore, the connection coefficients (and thus the normal fundamental forms) transform according to
\begin{align} \label{beta-transf}
\beta \mapsto m \beta m^{-1} + m \mathrm{d} \bigl(m^{-1}\bigr) ,
\end{align}
where adjacency is understood to imply matrix multiplication.
Note that for constant rotations, i.e., when $m$ is a constant matrix, $\beta$ transforms tensorially.

Because $\beta$ is not covariant under the group action, it is not a true invariant of the normal bundle (in that it does not transform linearly under the group action). Rather, its curvature 2-form
\begin{align} \label{norm-curv}
\mathcal{R} = \mathrm{d} \beta + [\beta, \beta]
\end{align}
does transform linearly under the group action and hence \textit{is} a true invariant of the normal bundle. In terms of the connection itself, the curvature 2-form, hereafter known as the \textit{normal curvature}, obeys the standard relation
\[
\mathcal{R}(\bar{u},\bar{v}) \xi = D_{\bar{u}} (D_{\bar{v}} \xi) - D_{\bar{v}} (D_{\bar{u}} \xi) - D_{[\bar{u},\bar{v}]} \xi
\]
for any $\bar{u},\bar{v} \in \Gamma(T \Lambda)$ and $\xi \in \Gamma(N \Lambda)$. The normal curvature is the final ingredient necessary to completely specify an embedded submanifold in $\mathbb{E}^d$ up to Euclidean motions.

To summarize, the three classical invariants of submanifold embeddings are $\bar{g}$, $\II$, and $\mathcal{R}$. It~is these tensors that appear in classical submanifold identities as described by Spivak~\cite{Spivak}, and given below for later use. First is the Gauss equation, relating the tangential component of the Riemann curvature tensor to that of the submanifold $\Lambda$,
\begin{align} \label{gauss-equation}
R_{abcd}^\top = \bar{R}_{abcd} + \II_{ad \alpha} \II_{bc \alpha} - \II_{bd \alpha} \II_{ac \alpha} ,
\end{align}
where $R^\top := \top (R)$, the projection of the Riemann curvature tensor to the submanifold. Similar notation will be used throughout.
Next is the Codazzi--Mainardi equation, relating another component of the Riemann curvature tensor to extrinsic invariants,
\begin{align} \label{codazzi}
R_{abc \alpha}^\top \= 2\bar{\nabla}_{[a} \II_{b]c \alpha} + 2\II_{c[b \beta} \beta_{a] \alpha \beta} .
\end{align}
Note that here, the right-hand side does not look explicitly tensorial -- however, a short calculation shows that it is indeed independent of the choice of orthonormal basis $\{n_{\alpha}\}$. Finally, we have the Ricci equation (which is trivial for hypersurfaces), relating another projection of the Riemann curvature to the normal curvature and the second fundamental form,
\[
R_{ab \alpha \beta}^\top = \mathcal{R}_{ab \alpha \beta} + \II^c_{a \beta} \II_{cb \alpha} - \II^c_{b \beta} \II_{ca \alpha} .
\]



% Original source lines 274--445
While the existence of a ``clean'' gauge choice can sometimes be taken for granted, uniqueness is much harder to demonstrate -- and indeed, it is not always present (see the discussion about the Gribov ambiguity above). Regardless, we can still study the extrinsic geometry of submanifold embeddings equipped with an orthonormal basis for the normal bundle. When a global orthonormal frame exists for $N \Lambda$, it is called a \textit{parallelization} of $N \Lambda$. Since we work locally, we can always assume that an orthonormal frame is specified for $N \Lambda$. We call such an embedding equipped with an orthonormal frame for the normal bundle a \textit{parallelized $($Riemannian$)$ submanifold embedding}, specified by the pair $(\Lambda \hookrightarrow (M,g), \{n_{\alpha}\})$.

\section[Defining maps for parallelized Riemannian submanifold embeddings]{Defining maps for parallelized Riemannian submanifold\\ embeddings} \label{unit-def-funcs}

Given a Riemannian submanifold embedding $\iota \colon \Lambda^{d-k} \hookrightarrow \bigl(M^d,g\bigr)$, we are interested in studying the embedding by examining the jets of a geometrically-adapted \textit{defining map} $s_{\alpha} \colon M \rightarrow \mathbb{R}^k$, which is a $k$-tuple of functions such that for any $p \in \Lambda$, we have that $s_{\alpha}(p) = 0$ for every $\alpha \in \{1, \ldots, k\}$ and ${\rm d}s_{1}(p) \wedge \cdots \wedge {\rm d}s_{k}(p) \neq 0$. That such families exist is well known, at least locally: for every $p \in \Lambda$, there exists a neighborhood $U \subset M$ such that $U \cap \Lambda$ is the zero locus of some defining map $s_{\alpha} \colon U \rightarrow \mathbb{R}^k$. However, there are many such defining maps, and in order for the jets of such a defining map to contain information about the embedding $\iota$, it must be determined canonically by the geometry. To realize this canonical defining map, we consider the \text{Gram matrix} $G_{\alpha \beta} := g^{-1}({\rm d}s_{\alpha},{\rm d}s_{\beta})$ of a given defining map $s_{\alpha}$. In the case of hypersurfaces, we canonically choose $G = \Id$, which is equivalent to the condition that $|{\rm d}s|^2 = 1$. Similarly, in the case of a higher-codimension submanifold we would like to achieve $G= \Id$. Because we will work locally, it is sufficient for our purposes to find a defining map $s_{\alpha}$ such that $G_{\alpha \beta} = \delta_{\alpha \beta} + \mathcal{O}(s^\infty)$, meaning we wish to find a defining map such that for any $m \in \mathbb{Z}_{\geq 1}$, we have that $G_{\alpha \beta} = \delta_{\alpha \beta} + \mathcal{O}(s^m)$. We call such a defining map a \textit{unit defining map}. Going forward, we will abuse notation by also denoting $g^{-1}({\rm d}s_{\alpha},\cdot)$ by $n_{\alpha}$. Note that, when clarification is necessary, we will decorate $n_{\alpha}$ with abstract indices (such as $n^a_{\alpha}$ and $n_{a \alpha}$) to indicate the whether we are referring to the vector or its metric dual.

\subsection{First order} \label{fo-Riem}
Given the parallelization $\{n_{\alpha}\}$ of $N \Lambda$ and any defining map $\hat{s}_{\alpha}$ for $\Lambda$, there exists a linear combination of components of $\hat{s}_{\alpha}$ denoted by $s_{\alpha}$ such that $g^{-1}({\rm d}s_{\alpha},\cdot) \= n_{\alpha}$. This linear combination provides a unique $s_{\alpha}$ modulo terms that are at least second order in a defining function. The behavior of ${\rm d}s_{\alpha}$ away from $\Lambda$, however, is harder to control. Our goal is to obtain a canonical family of unit defining functions adapted to our parallelization satisfying $g^{-1}({\rm d}s_{\alpha},{\rm d}s_{\beta}) = \delta_{\alpha \beta}$. However, this will not be possible in general. Instead, we will choose $s_{\alpha}$ to, in some sense, get as close to the identity Gram matrix condition as possible.
We begin with an iterative procedure. Because the parallelization of $N \Lambda$ is defined to be orthonormal, we can certainly obtain $G_{\alpha \beta} \= \delta_{\alpha \beta}$, so we assume this going forward. This constraint defines a unique family of defining maps that are aligned to the parallelization of $N \Lambda$ and whose elements have a difference given~by \smash{$s'_{\alpha} - s_{\alpha} = \mathcal{O}\bigl(s^2\bigr)$}. For any of these defining maps, it thus follows that
\begin{align} \label{0-order}
G_{\alpha \beta} = \delta_{\alpha \beta} + F^{(1)}_{\alpha \beta \gamma} s_{\gamma} ,
\end{align}
where $F^{(1)}_{\alpha \beta \gamma}$ represents an array of smooth functions on $M$. Our goal will be to pare down this family of defining maps by controlling \smash{$F^{(1)}$}.

However, note that equation~(\ref{0-order}) does not uniquely specify \smash{$F^{(1)}_{\alpha \beta \gamma}$} as an array of functions on $C^\infty M$. To see this, consider the mapping
\[
F^{(1)}_{\alpha \beta \gamma} \mapsto F^{(1)}_{\alpha \beta \gamma} + F^{(1,1)}_{\alpha \beta \gamma \delta} s_{\delta} ,
\]
where \smash{$F^{(1,1)}$} is any array of functions so that \smash{$F^{(1,1)}_{\alpha \beta \gamma \delta} = F^{(1,1)}_{\alpha \beta [\gamma \delta]}$}. This mapping clearly preserves equation~(\ref{0-order}), and so we see that \smash{$F^{(1)}$} is only uniquely determined along $\Lambda$. Nonetheless, we \textit{may} choose \smash{$F^{(1)}$} in a collar neighborhood of $\Lambda$ in a non-canonical way. Going forward, we will assume such a choice has been made; later, we will show that this choice actually yields a unique result.

Because of the symmetry of $G_{\alpha \beta}$, we have that \smash{$F^{(1)}_{\alpha \beta \gamma} = F^{(1)}_{(\alpha \beta) \gamma}$}, and thus as a representation of $S_k$, we can see that the tensor structure of \smash{$F^{(1)}$} in the normal bundle takes the form
\[
\raisebox{3pt}{\scalebox{0.4}{\ydiagram{2,1}}} \; \oplus \; \scalebox{0.4}{\ydiagram{3}} .
\]
For what follows, we will abuse notation and write, for example, that $F^{(1)} \in C^\infty (M,\raisebox{3pt}{\scalebox{0.2}{\ydiagram{2,1}}} \; \oplus \; \raisebox{2pt}{\scalebox{0.2}{\ydiagram{3}}})$.

Now observe that under the mapping
\[
s_{\alpha} \mapsto \tilde{s}_{\alpha} := s_{\alpha} + A^{(1)}_{\alpha \gamma_1 \gamma_2} s_{\gamma_1} s_{\gamma_2} ,
\]
where $A^{(1)}$ is a list of smooth functions on $M$, we have that the zeroth order part of the Gram matrix condition given in equation~(\ref{0-order}) is preserved while changing $F^{(1)}$. Indeed, by inspection we may consider only those \smash{$A^{(1)}$} satisfying \smash{$A^{(1)}_{\alpha \gamma_1 \gamma_2} = A^{(1)}_{\alpha (\gamma_1 \gamma_2)}$}, and so $A^{(1)} \in C^\infty (M , \raisebox{3pt}{\scalebox{0.2}{\ydiagram{2,1}}} \; \oplus \; \raisebox{2pt}{\scalebox{0.2}{\ydiagram{3}}})$. Evidently, it may be possible to choose \smash{$A^{(1)}$} so that it cancels \smash{$F^{(1)}$}.

Differentiating the above display and using the definition of the Gram matrix for the new defining map (combined with equation~\eqref{0-order}), we obtain
\begin{align}
\tilde{G}_{\alpha \beta} &{}= \delta_{\alpha \beta}+ \bigl(F^{(1)}_{\alpha \beta \gamma} + 4 A^{(1)}_{(\alpha \beta) \gamma}\bigr) s_{\gamma} \nonumber
\\
&\quad{}+ \bigl(4 A^{(1)}_{\alpha (\rho \gamma_1)} A^{(1)}_{\beta (\rho \gamma_2)} + 2 n_{a (\alpha} \nabla^a A^{(1)}_{\beta) (\gamma_1 \gamma_2)} + 4 A^{(1)}_{\alpha (\rho \gamma_1)} F^{(1)}_{\beta \rho \gamma_2} \bigr) s_{\gamma_1} s_{\gamma_2}
\nonumber\\
&\quad{}+ \bigl(4A^{(1)}_{\alpha (\rho \gamma_1)} \nabla_{n_{\rho}} A^{(1)}_{\beta \gamma_2 \gamma_3} + 4 A^{(1)}_{\alpha (\omega_1 \gamma_2)} A^{(1)}_{\beta (\omega_2 \gamma_2)} F^{(1)}_{\omega_1 \omega_2 \gamma_3} \bigr) s_{\gamma_1} s_{\gamma_2} s_{\gamma_3} \nonumber\\
&\quad{}+ \bigl(\nabla^a A^{(1)}_{\alpha (\gamma_1 \gamma_2)}\bigr) \bigl(\nabla_a A^{(1)}_{\beta (\gamma_3 \gamma_4)} \bigr) s_{\gamma_1} s_{\gamma_2} s_{\gamma_3} s_{\gamma_4} ,
\label{1-order-correction}
\end{align}
where the above is symmetric in $\alpha \beta$ (even when the round brackets are not displayed).
Observe that in order to fix $A^{(1)}$, we must express it in terms of $F^{(1)}$. To that end, we fix an ansatz: \smash{$A^{(1)}_{\alpha \beta \gamma} = a F^{(1)}_{\alpha \beta \gamma} + b F^{(1)}_{\gamma \alpha \beta} + c F^{(1)}_{\beta \gamma \alpha}$} and then solve explicitly; these are the only three terms~possible~as \smash{$F^{(1)}_{\alpha \beta \gamma} = F^{(1)}_{(\alpha \beta) \gamma}$}. We find that
\[
A^{(1)}_{\alpha \beta \gamma} = -\frac{1}{4}\bigl(F^{(1)}_{\alpha \beta \gamma} + F^{(1)}_{\gamma \alpha \beta} - F^{(1)}_{\beta \gamma \alpha}\bigr)
\]
sets the first order term in equation~\eqref{1-order-correction} to zero identically, and thus we have found $s_{\alpha}$ such that
\[
G_{\alpha \beta} = \delta_{\alpha \beta} + \mathcal{O}\bigl(s^2\bigr) ,
\]
fixing $s_{\alpha}$ to second order in a way dependent on the original choice of extension of $F^{(1)}$.


With this result, we find a useful identity for the normal fundamental form $\beta$ associated with the parallelization of $N \Lambda$:
\begin{align}
n^b_{\alpha} \nabla_a n_{b \beta} \={}& n^b_{[\alpha} \nabla_a n_{b \beta]} + n^b_{(\alpha} \nabla_a n_{b \beta)}
\= n^b_{[\alpha} \bigl(\nabla_a^\top + n_{a \gamma} n^c_{\gamma} \nabla_c \bigr) n_{b \beta]} + \frac{1}{2} \nabla_a n^b_{(\alpha} n_{b \beta)} \nonumber \\
\={}& \beta_{a \alpha \beta} + n_{a \gamma} n^c_{\gamma} n^b_{[\alpha} \nabla_c n_{b \beta]}
\= \beta_{a \alpha \beta} + n_{a \gamma} n^c_{\gamma} n^b_{[\alpha} \nabla_b n_{c \beta]} \nonumber \\
\={}& \beta_{a \alpha \beta} - n_{a \gamma} n^b_{[\alpha} n^c_{\beta]} \nabla_b n_{c \gamma}
\= \beta_{a \alpha \beta} . \label{ndn-Riem}
\end{align}
In the above computation, we used that when $n_{\alpha} = {\rm d} s_{\alpha}$, we evidently have that $\nabla_{[a} n_{b]\beta} = 0$. Observe from this calculation that \smash{$[n_{\alpha},n_{\beta}]^a \= 2 \beta^a_{\alpha \beta}$}. Further, it thus follows that
\begin{align} \label{nabla-n}
\nabla_a n_{b \beta} \= \II_{ab \beta} + n_{b \alpha} \beta_{a \alpha \beta} + n_{a \alpha} \beta_{b \alpha \beta} .
\end{align}


\subsection{Second order} %\label{2o-Riem}
Section~\ref{fo-Riem} established that for a given parallelized Riemannian submanifold embedding, there exists a unique family of defining maps related by $s'_{\alpha} - s_{\alpha} = \mathcal{O}\bigl(s^3\bigr)$ such that
\begin{align} \label{1-order}
G_{\alpha \beta} = \delta_{\alpha \beta} + F^{(2)}_{\alpha \beta \gamma_1 \gamma_2} s_{\gamma_1} s_{\gamma_2} .
\end{align}
As in that section, our goal is to pare down this family by controlling $F^{(2)}$. From equation~(\ref{1-order-correction}), $F^{(2)}$ is uniquely determined in a collar neighborhood of $\Lambda$ given the initial choice of $F^{(1)}$. Our goal will be to adjust $s_{\alpha}$ so that we may remove as much of $F^{(2)}$ as possible.

Using the same reasoning as in Section~\ref{fo-Riem}, we have that $F^{(2)} \in C^\infty (M, \raisebox{3pt}{\scalebox{0.2}{\ydiagram{2,2}}} \; \oplus \raisebox{3pt}{\scalebox{0.2}{\ydiagram{3,1}}} \; \oplus \; \raisebox{1.5pt}{\scalebox{0.2}{\ydiagram{4}}})$. So, we attempt to fix the family of defining functions to cancel $F^{(2)}$ by applying the map
\[
s_{\alpha} \mapsto \tilde{s}_{\alpha} := s_{\alpha} + A^{(2)}_{\alpha \gamma_1 \gamma_2 \gamma_3} s_{\gamma_1} s_{\gamma_2} s_{\gamma_3} ,
\]
which clearly preserves the structure of the Gram matrix condition~(\ref{1-order}) while possibly changing~$F^{(2)}$. However, note that $A^{(2)} \in C^\infty (M, \raisebox{3pt}{\scalebox{0.2}{\ydiagram{3,1}}} \; \oplus \; \raisebox{1.5pt}{\scalebox{0.2}{\ydiagram{4}}})$. Because $A^{(2)}$ and $F^{(2)}$ can live in different representations of $S_k$, it is not always possible to construct $A^{(2)}$ such that $F^{(2)} \mapsto 0$. Nonetheless, one can always find $A^{(2)}$ so that
\[
F^{(2)}_{\alpha \beta \gamma_1 \gamma_2} \mapsto \tilde{F}^{(2)}_{(\alpha \beta)(\gamma_1 \gamma_2)} \in C^\infty (M,\raisebox{3pt}{\scalebox{0.2}{\ydiagram{2,2}}}).
\]



From the above construction, for a given parallelization of $N \Lambda$, there always exists at least one defining map such that
\begin{align*} %\label{2-obs}
G_{\alpha \beta} = \delta_{\alpha \beta} + F^{(2)}_{\alpha \beta \gamma_1 \gamma_2} s_{\gamma_1} s_{\gamma_2} ,
\end{align*}
where $F^{(2)} \in C^\infty (M,\raisebox{3pt}{\scalebox{0.2}{\ydiagram{2,2}}})$ is entirely determined by the parallelization of $N \Lambda$ and the choice of~$F^{(1)}$. We call defining maps that are in this infinite family \textit{associated defining maps}. Along~$\Lambda$, $F^{(2)}|_{\Lambda}$ is a first obstruction to producing a holonomic basis for the extension of the normal bundle. In fact, $F^{(2)}|_{\Lambda}$ can be computed in terms of tensors specified by the parallelization. We now carry out this computation, and going forward we will work with the family of associated defining maps to the parallelization.


Now, observe that
\[
P_{\scalebox{0.2}{\ydiagram{2,2}}}\, n^a_{\gamma_1} n^b_{\gamma_2} \nabla_a \nabla_b G_{\alpha \beta} \= 2 F^{(2)}_{\alpha \beta \gamma_1 \gamma_2} ,
\]
where $P_{\scalebox{0.2}{\ydiagram{2,2}}}$ is the projector onto the $\raisebox{3pt}{{\scalebox{0.2}{\ydiagram{2,2}}}}$ component of a rank-4 tensor in $S^k$. In calculations that will follow, it is useful to have explicit formulas for relevant projection operators. To that end, we find that
\begin{align*}
P_{\scalebox{0.2}{\ydiagram{2,2}}}\, F_{\alpha \beta \gamma \delta} = \frac{1}{12} \bigl(&F_{\alpha \beta \gamma \delta} - F_{\gamma \beta \alpha \delta} - F_{\alpha \delta \gamma \beta} + F_{\gamma \delta \alpha \beta} + F_{\beta \alpha \gamma \delta} - F_{\gamma \alpha \beta \delta} - F_{\beta \delta \gamma \alpha} + F_{\gamma \delta \beta \alpha} \\
&+ F_{\alpha \beta \delta \gamma}- F_{\delta \beta \alpha \gamma} - F_{\alpha \gamma \delta \beta} + F_{\delta \gamma \alpha \beta} + F_{\beta \alpha \delta \gamma} - F_{\delta \alpha \beta \gamma} - F_{\beta \gamma \delta \alpha} + F_{\delta \gamma \beta \alpha}\bigr) ,
\end{align*}
and
\begin{align*}
P_{\scalebox{0.2}{\ydiagram{3,1}}}\, F_{\alpha \beta \gamma \delta} = \frac{1}{8}\bigl(&F_{\alpha \beta \gamma \delta} - F_{\delta \beta \gamma \alpha} + F_{\beta \alpha \gamma \delta} - F_{\delta \alpha \gamma \beta} + F_{\gamma \beta \alpha \delta} - F_{\delta \beta \alpha \gamma} + F_{\alpha \gamma \beta \delta} - F_{\delta \gamma \beta \alpha}
\\&+ F_{\beta \gamma \alpha \delta}- F_{\delta \gamma \alpha \beta} + F_{\gamma \alpha \beta \delta} - F_{\delta \alpha \beta \gamma} \bigr) .
\end{align*}
We can now compute directly
\begin{align*}
P_{\scalebox{0.2}{\ydiagram{2,2}}}\, n^a_{\gamma_1} n^b_{\gamma_2} \nabla_a \nabla_b n^c_{\alpha} n_{c \beta} ={}& P_{\scalebox{0.2}{\ydiagram{2,2}}}\, n^{(a}_{\gamma_1} n^{b)}_{\gamma_2} \bigl[ 2 (\nabla n_{\alpha})^c_{(a} (\nabla n_{\beta})_{b) c} + 2 n^c_{(\alpha|} \nabla_a \nabla_b n_{c| \beta)} \bigr] \\
\={}& -2 \beta_{c \alpha (\gamma_1 } \beta^c_{\gamma_2) \beta} - \frac{2}{3} R_{\gamma_1 (\alpha \beta) \gamma_2} .
\end{align*}
Putting this together, we have that
\begin{align} \label{Riem-2o-obstruction}
F^{(2)}_{\alpha \beta \gamma_1 \gamma_2} \= - \beta_{c \alpha (\gamma_1 } \beta^c_{\gamma_2) \beta} - \frac{1}{3} R_{\gamma_1 (\alpha \beta) \gamma_2} .
\end{align}
It is important to note that $F^{(2)}$ is (as yet) dependent on the choice of extension of $F^{(1)}|_{\Lambda}$. However, what this computation verifies is that this choice of extension yields a unique obstruction $F^{(2)}|_{\Lambda}$ that depends only on the parallelization of $N \Lambda$ and not on some choice of how the parallelization or $F^{(1)}|_{\Lambda}$ extends off $\Lambda$.

Thus, the obstruction in equation~(\ref{Riem-2o-obstruction}) is a true invariant of the parallelized Riemannian submanifold embedding. As such, it may be of independent interest to study those structures for which this obstruction vanishes. While there is indeed an obstruction, we may proceed regardless.


\subsection{Higher orders} \label{ho-Riem}
We now seek to generalize the result of the previous subsection to all orders. As noted, at second order $F^{(2)}|_{\Lambda}$ is unique, despite the choice made of extension of $F^{(1)}$ off $\Lambda$. We would like to show that this phenomenon continues. This result is captured by the following theorem.

\begin{Theorem} \label{general-ho-ext}
Let $\Lambda^{d-k} \hookrightarrow \bigl(M^d,g\bigr)$ be parallelized by $\{n_{\alpha}\}$. Then, in a collar neighborhood of~$\Lambda$, there exists formally a unique defining map determined by $\{n_{\alpha}\}$ such that
\begin{align} \label{window-expansion}
G_{\alpha \beta} = \delta_{\alpha \beta} + F^{(2)}_{\alpha \beta \gamma_1 \gamma_2} s_{\gamma_1} s_{\gamma_2} + F^{(3)}_{\alpha \beta \gamma_1 \gamma_2 \gamma_3} s_{\gamma_1} s_{\gamma_2} s_{\gamma_3} + \cdots,
\end{align}
where each $F^{(i)}$ is canonically determined to all orders and lives in a generalized window tensor representation of $S^k$.
\end{Theorem}
For clarity, a generalized window tensor representation of $S^k$ corresponds to a Young diagram with a first row of length at least two, and second row of length two.
\begin{proof}
We first show that one may find a defining map that satisfies equation~(\ref{window-expansion}) via induction. To do so, we begin by fixing an arbitrary extension of $F^{(1)}|_{\Lambda}$ as described in Section~\ref{fo-Riem}. Given such an extension, the base case is handled by the previous subsection, so that there exists a defining map so that \smash{$G = \delta + \mathcal{O}\bigl(s^2\bigr)$}. Now, suppose that there exists a defining map such that
\[
G_{\alpha \beta} = \delta_{\alpha \beta} + F^{(2)}_{\alpha \beta \gamma_1 \gamma_2} s_{\gamma_1} s_{\gamma_2} + \cdots + F^{(m)}_{\alpha \beta \gamma_1 \cdots \gamma_m} s_{\gamma_1} \cdots s_{\gamma_m} + F^{(m+1)}_{\alpha \beta \gamma_1 \cdots \gamma_{m+1}} s_{\gamma_1} \cdots s_{\gamma_{m+1}} ,
\]
where for each $2 \leq i \leq m$, $F^{(i)}$ occupies a generalized window tensor representation. Now let
\[
s_{\alpha} \mapsto \tilde{s}_{\alpha} := s_{\alpha} + A^{(m+1)}_{\alpha \gamma_1 \cdots \gamma_{m+2}} s_{\gamma_1} \cdots s_{\gamma_{m+2}} .
\]
Then
\[
\tilde{G}_{\alpha \beta} = G_{\alpha \beta} + 2(m+1) A^{(m+1)}_{(\alpha \beta) \gamma_1 \cdots \gamma_{m+1}} s_{\gamma_1} \cdots s_{\gamma_{m+1}} + \mathcal{O}\bigl(s^{m+2}\bigr) .
\]
From the representation theory of the normal bundle, we find that $A^{(m+1)}$ may thus be chosen in a collar neighborhood of $\Lambda$ so that $F^{(m+1)}$ occupies the generalized window tensor representation. This completes the induction, and thus given an arbitrary extension of $F^{(1)}|_{\Lambda}$ off $\Lambda$, we have uniquely determined a defining map $s_{\alpha}$ that satisfies equation~(\ref{window-expansion}).

Next, we must show that this defining map is independent of the choice of $F^{(1)}$. Let $F^{(1)}$ and $F^{(1)'}$ be distinct extensions of $F^{(1)}|_{\Lambda}$. Then we may produce unique defining maps $s$ and $s'$ respectively associated to each of these choices as described above. Now because the differentials of the defining maps necessarily agree on the submanifold, we must have that
\[
s'_{\alpha} = s_{\alpha} + M_{\alpha \gamma_1 \gamma_2} s_{\gamma_1} s_{\gamma_2}
\]
for some $M_{\alpha \gamma_1 \gamma_2} \in C^\infty M \times (\raisebox{3pt}{\scalebox{0.2}{\ydiagram{2,1}}} \; \oplus \; \raisebox{2pt}{\scalebox{0.2}{\ydiagram{3}}})$. We may now compare the Gram matrices,
\[
G'_{\alpha \beta} = G_{\alpha \beta} + 2 M_{(\alpha \beta) \gamma_1} s_{\gamma_1} + \mathcal{O}\bigl(s^2\bigr) .
\]
However, observe that $G'_{\alpha \beta}$ and $G_{\alpha \beta}$ must be the identity to second order, and so~\smash{$M_{(\alpha \beta) \gamma} \= 0$}. So because \smash{$M_{\alpha \beta \gamma} \= M_{[\alpha \beta] \gamma}$},
\[
M_{\alpha \beta \gamma} \= -M_{\beta \alpha \gamma} \= -M_{\beta \gamma \alpha} \= M_{\gamma \beta \alpha} \= M_{\gamma \alpha \beta} \= -M_{\alpha \gamma \beta} \= -M_{\alpha \beta \gamma} .
\]
Thus, we have that $M_{\alpha \beta \gamma} \= 0$.

So this means that
\[
s'_{\alpha} = s_{\alpha} + M_{\alpha \gamma_1 \gamma_2 \gamma_3} s_{\gamma_1} s_{\gamma_2} s_{\gamma_3} .
\]
By the same considerations, we have that
\[
G'_{\alpha \beta} - G_{\alpha \beta} = 4 M_{(\alpha \beta) \gamma_1 \gamma_2} s_{\gamma_1} s_{\gamma_2} .
\]
However, at second order the left-hand side necessarily occupies a window tensor representation~-- but $M$ is symmetric in three indices, and thus cannot contribute to this tensor representation. Therefore, it must vanish along $\Lambda$, and hence $s' - s = \mathcal{O}\bigl(s^4\bigr)$. By a similar argument, we may apply induction to find that $s'_{\alpha} = s_{\alpha}$ to arbitrarily high order. Thus, the defining map is independent of the choice of the extension of $F^{(1)}|_{\Lambda}$.
\end{proof}
